\documentclass[11pt,letterpaper]{article}
\usepackage[utf8]{inputenc}
\usepackage{amsmath}
\usepackage{amssymb}
\usepackage{geometry}
\geometry{margin=1in}
\usepackage{booktabs}
\usepackage{hyperref}
\usepackage{enumitem}
\usepackage{xcolor}

\title{\textbf{Chapter 4 (Continued): Higher-Order Linear Differential Equations}\\
\large $n$-th Order Homogeneous ODEs, IVPs, Equation Reconstruction, and Complex Root Logic\\
\normalsize Ecuaciones Diferenciales Lineales de Orden $n$}
\author{Gemini Notebook}
\date{\today}

\begin{document}
\maketitle

\section{Theoretical Overview / Resumen Teórico}

Higher-order homogeneous linear differential equations with constant real coefficients take the general form:
\begin{equation}
    a_n y^{(n)} + a_{n-1} y^{(n-1)} + \dots + a_1 y' + a_0 y = 0, \quad a_n \neq 0, \quad a_i \in \mathbb{R}
\end{equation}

\subsection{Characteristic Polynomial and Fundamental Set of Solutions}
Assuming a solution of the form $y = e^{rx}$, we obtain the $n$-th degree auxiliary / characteristic polynomial:
\begin{equation}
    P(r) = a_n r^n + a_{n-1} r^{n-1} + \dots + a_1 r + a_0 = 0
\end{equation}
By the Fundamental Theorem of Algebra, $P(r)$ has exactly $n$ complex roots (counting multiplicity). The structure of the homogeneous general solution $y_h(x) = c_1 y_1(x) + c_2 y_2(x) + \dots + c_n y_n(x)$ is determined as follows:

\begin{enumerate}
    \item \textbf{Real Distinct Roots:} Each simple real root $r_k$ contributes $e^{r_k x}$.
    \item \textbf{Real Repeated Roots:} A real root $r_k$ of multiplicity $m$ contributes $m$ linearly independent solutions:
    \[ e^{r_k x}, \, x e^{r_k x}, \, x^2 e^{r_k x}, \, \dots, \, x^{m-1} e^{r_k x} \]
    \item \textbf{Complex Conjugate Roots:} Since coefficients $a_i$ are real, non-real roots occur in conjugate pairs $\alpha \pm i \beta$. A simple pair contributes:
    \[ e^{\alpha x} \cos(\beta x) \quad \text{and} \quad e^{\alpha x} \sin(\beta x) \]
    \item \textbf{Repeated Complex Conjugate Roots:} A pair $\alpha \pm i \beta$ of multiplicity $m$ contributes $2m$ solutions:
    \[ e^{\alpha x} \cos(\beta x), \, x e^{\alpha x} \cos(\beta x), \, \dots, \, x^{m-1} e^{\alpha x} \cos(\beta x) \]
    \[ e^{\alpha x} \sin(\beta x), \, x e^{\alpha x} \sin(\beta x), \, \dots, \, x^{m-1} e^{\alpha x} \sin(\beta x) \]
\end{enumerate}

\subsection{Reconstructing the Differential Equation from Solutions}
If a fundamental set of solutions $\{y_1, y_2, \dots, y_n\}$ is known, we can identify the corresponding characteristic roots $r_1, r_2, \dots, r_n$ and build the characteristic polynomial:
\begin{equation}
    P(r) = (r - r_1)(r - r_2) \cdots (r - r_n) = 0
\end{equation}
Expanding $P(r)$ into $a_n r^n + a_{n-1} r^{n-1} + \dots + a_0 = 0$ yields the differential equation operator $P(D)y = 0$, where $D^k y = y^{(k)}$.

\newpage
\section{Problem Set / Banco de Ejercicios}

\subsection*{Part A: 10 Exercises - Finding the Homogeneous Solution $y_h$}
Find the general solution $y_h(x)$ for each of the following $n$-th order homogeneous linear differential equations.

\begin{align}
    & y''' - 6y'' + 11y' - 6y = 0 \label{eq:nth_1} \\
    & y''' - 3y'' + 3y' - y = 0 \label{eq:nth_2} \\
    & y''' + y' = 0 \label{eq:nth_3} \\
    & y^{(4)} - 16y = 0 \label{eq:nth_4} \\
    & y^{(4)} + 8y'' + 16y = 0 \label{eq:nth_5} \\
    & y''' - y'' - y' + y = 0 \label{eq:nth_6} \\
    & y^{(4)} - 5y'' + 4y = 0 \label{eq:nth_7} \\
    & y''' + 2y'' + 4y' + 8y = 0 \label{eq:nth_8} \\
    & y^{(5)} - 2y^{(4)} + y''' = 0 \label{eq:nth_9} \\
    & y^{(4)} + 2y'' + y = 0 \label{eq:nth_10}
\end{align}

\subsection*{Part B: 5 Exercises - Initial Value Problems (IVPs)}
Solve the following higher-order initial value problems.

\begin{align}
    & y''' - 2y'' - y' + 2y = 0, \quad y(0) = 1, \, y'(0) = 2, \, y''(0) = 0 \label{eq:ivp_1} \\
    & y''' + 4y' = 0, \quad y(0) = 0, \, y'(0) = 4, \, y''(0) = 8 \label{eq:ivp_2} \\
    & y''' - 3y'' + 3y' - y = 0, \quad y(0) = 1, \, y'(0) = 2, \, y''(0) = 3 \label{eq:ivp_3} \\
    & y^{(4)} - y = 0, \quad y(0) = 1, \, y'(0) = 0, \, y''(0) = 1, \, y'''(0) = 0 \label{eq:ivp_4} \\
    & y''' + y'' + y' + y = 0, \quad y(0) = 2, \, y'(0) = -1, \, y''(0) = 0 \label{eq:ivp_5}
\end{align}

\subsection*{Part C: 5 Exercises - Reconstructing the Differential Equation}
For each given fundamental set of solutions, reconstruct the minimal-order homogeneous linear differential equation with constant real coefficients $P(D)y = 0$.

\begin{enumerate}[label=\textbf{\arabic*.}]
    \item $S = \{1, \, e^{2x}, \, e^{-2x}\}$
    \item $S = \{e^x, \, x e^x, \, x^2 e^x\}$
    \item $S = \{\cos(3x), \, \sin(3x), \, e^x\}$
    \item $S = \{1, \, x, \, e^{4x}\}$
    \item $S = \{e^{-x}, \, e^x \cos(2x), \, e^x \sin(2x)\}$
\end{enumerate}

\subsection*{Part D: 3 Logic / Clue Problems - Complex Roots and Minimal Order}
Apply the theory of real-coefficient differential equations (Dennis Zill style) to solve these conceptual puzzles.

\begin{enumerate}[label=\textbf{Problem \arabic*:}, leftmargin=*]
    \item \textbf{Complex Solution Clue:} A homogeneous linear differential equation with \textit{constant real coefficients} has $y_1(x) = e^{3x} \sin(2x)$ as one of its solutions.
    \begin{enumerate}[label=(\alph*)]
        \item What is the minimal order that this differential equation can have?
        \item State all roots of the minimal characteristic polynomial and write the differential equation.
        \item If it is additionally known that $y_2(x) = 5$ is also a solution, what is the new minimal order and the corresponding differential equation?
    \end{enumerate}

    \item \textbf{Repeated Complex Root Clue:} A homogeneous linear ODE with real constant coefficients has $y_1(x) = x e^{-x} \cos(x)$ as a solution.
    \begin{enumerate}[label=(\alph*)]
        \item Identify all roots of the characteristic equation implied by this solution, including their multiplicities.
        \item Determine the minimal order $n$ of the differential equation.
        \item Find the explicit differential equation in operator notation $P(D)y = 0$.
    \end{enumerate}

    \item \textbf{Reconstruction from Partial Solutions and IVP:} 
    A 3rd-order homogeneous linear differential equation with constant real coefficients has solutions $y_1(x) = \sin(x)$ and $y_2(x) = e^{2x}$.
    \begin{enumerate}[label=(\alph*)]
        \item Deduce the full set of 3 characteristic roots and write the differential equation.
        \item Solve the initial value problem with conditions $y(0) = 0$, $y'(0) = 3$, and $y''(0) = 5$.
    \end{enumerate}
\end{enumerate}

\newpage
\section{Solutions and Step-by-Step Manual / Solucionario}

\subsection*{Part A: Homogeneous Solutions $y_h$}

\subsubsection*{Solution 1: $y''' - 6y'' + 11y' - 6y = 0$}
Characteristic equation: $r^3 - 6r^2 + 11r - 6 = 0$.
Testing $r = 1$: $1 - 6 + 11 - 6 = 0$. So $(r-1)$ is a factor.
Factoring: $(r-1)(r^2 - 5r + 6) = (r-1)(r-2)(r-3) = 0$.
Roots: $r_1 = 1, r_2 = 2, r_3 = 3$.
\[ y_h(x) = c_1 e^x + c_2 e^{2x} + c_3 e^{3x} \]

\subsubsection*{Solution 2: $y''' - 3y'' + 3y' - y = 0$}
Characteristic equation: $r^3 - 3r^2 + 3r - 1 = (r-1)^3 = 0$.
Root: $r = 1$ with multiplicity 3.
\[ y_h(x) = c_1 e^x + c_2 x e^x + c_3 x^2 e^x = (c_1 + c_2 x + c_3 x^2) e^x \]

\subsubsection*{Solution 3: $y''' + y' = 0$}
Characteristic equation: $r^3 + r = r(r^2 + 1) = 0$.
Roots: $r_1 = 0, r_{2,3} = \pm i$ ($\alpha = 0, \beta = 1$).
\[ y_h(x) = c_1 e^{0x} + c_2 \cos(x) + c_3 \sin(x) = c_1 + c_2 \cos(x) + c_3 \sin(x) \]

\subsubsection*{Solution 4: $y^{(4)} - 16y = 0$}
Characteristic equation: $r^4 - 16 = (r^2 - 4)(r^2 + 4) = (r-2)(r+2)(r^2 + 4) = 0$.
Roots: $r_{1,2} = \pm 2$, $r_{3,4} = \pm 2i$.
\[ y_h(x) = c_1 e^{2x} + c_2 e^{-2x} + c_3 \cos(2x) + c_4 \sin(2x) \]

\subsubsection*{Solution 5: $y^{(4)} + 8y'' + 16y = 0$}
Characteristic equation: $r^4 + 8r^2 + 16 = (r^2 + 4)^2 = 0$.
Roots: $r = \pm 2i$ with multiplicity 2 each.
\[ y_h(x) = (c_1 + c_2 x) \cos(2x) + (c_3 + c_4 x) \sin(2x) \]

\subsubsection*{Solution 6: $y''' - y'' - y' + y = 0$}
Characteristic equation: $r^3 - r^2 - r + 1 = r^2(r-1) - (r-1) = (r^2 - 1)(r-1) = (r-1)^2(r+1) = 0$.
Roots: $r_1 = r_2 = 1$ (multiplicity 2), $r_3 = -1$.
\[ y_h(x) = c_1 e^x + c_2 x e^x + c_3 e^{-x} \]

\subsubsection*{Solution 7: $y^{(4)} - 5y'' + 4y = 0$}
Characteristic equation: $r^4 - 5r^2 + 4 = (r^2 - 1)(r^2 - 4) = (r-1)(r+1)(r-2)(r+2) = 0$.
Roots: $r = \pm 1, \pm 2$.
\[ y_h(x) = c_1 e^x + c_2 e^{-x} + c_3 e^{2x} + c_4 e^{-2x} \]

\subsubsection*{Solution 8: $y''' + 2y'' + 4y' + 8y = 0$}
Characteristic equation: $r^2(r+2) + 4(r+2) = (r+2)(r^2+4) = 0$.
Roots: $r_1 = -2, r_{2,3} = \pm 2i$.
\[ y_h(x) = c_1 e^{-2x} + c_2 \cos(2x) + c_3 \sin(2x) \]

\subsubsection*{Solution 9: $y^{(5)} - 2y^{(4)} + y''' = 0$}
Characteristic equation: $r^3(r^2 - 2r + 1) = r^3(r-1)^2 = 0$.
Roots: $r = 0$ (multiplicity 3), $r = 1$ (multiplicity 2).
\[ y_h(x) = c_1 + c_2 x + c_3 x^2 + c_4 e^x + c_5 x e^x \]

\subsubsection*{Solution 10: $y^{(4)} + 2y'' + y = 0$}
Characteristic equation: $r^4 + 2r^2 + 1 = (r^2 + 1)^2 = 0$.
Roots: $r = \pm i$ with multiplicity 2 each.
\[ y_h(x) = (c_1 + c_2 x) \cos(x) + (c_3 + c_4 x) \sin(x) \]

\subsection*{Part B: Initial Value Problems (IVPs)}

\subsubsection*{Solution 11: $y''' - 2y'' - y' + 2y = 0, \, y(0)=1, \, y'(0)=2, \, y''(0)=0$}
Characteristic equation: $r^2(r-2) - (r-2) = (r-2)(r^2-1) = (r-2)(r-1)(r+1) = 0$.
General solution: $y(x) = c_1 e^x + c_2 e^{-x} + c_3 e^{2x}$.
Derivatives:
\[ y'(x) = c_1 e^x - c_2 e^{-x} + 2c_3 e^{2x} \]
\[ y''(x) = c_1 e^x + c_2 e^{-x} + 4c_3 e^{2x} \]
Applying initial conditions at $x=0$:
\[ c_1 + c_2 + c_3 = 1 \]
\[ c_1 - c_2 + 2c_3 = 2 \]
\[ c_1 + c_2 + 4c_3 = 0 \]
Subtracting equation 1 from equation 3: $3c_3 = -1 \implies c_3 = -1/3$.
Then $c_1 + c_2 = 1 - c_3 = 4/3$, and $c_1 - c_2 = 2 - 2c_3 = 8/3$.
Adding: $2c_1 = 12/3 = 4 \implies c_1 = 2$.
$c_2 = 4/3 - 2 = -2/3$.
Particular Solution:
\[ y(x) = 2e^x - \frac{2}{3}e^{-x} - \frac{1}{3}e^{2x} \]

\subsubsection*{Solution 12: $y''' + 4y' = 0, \, y(0)=0, \, y'(0)=4, \, y''(0)=8$}
Roots: $r = 0, \pm 2i$.
General solution: $y(x) = c_1 + c_2 \cos(2x) + c_3 \sin(2x)$.
Derivatives:
\[ y'(x) = -2c_2 \sin(2x) + 2c_3 \cos(2x) \]
\[ y''(x) = -4c_2 \cos(2x) - 4c_3 \sin(2x) \]
Applying ICs:
\[ y(0) = c_1 + c_2 = 0 \implies c_1 = -c_2 \]
\[ y'(0) = 2c_3 = 4 \implies c_3 = 2 \]
\[ y''(0) = -4c_2 = 8 \implies c_2 = -2 \implies c_1 = 2 \]
Particular Solution:
\[ y(x) = 2 - 2\cos(2x) + 2\sin(2x) \]

\subsubsection*{Solution 13: $y''' - 3y'' + 3y' - y = 0, \, y(0)=1, \, y'(0)=2, \, y''(0)=3$}
Roots: $r = 1$ (multiplicity 3).
General solution: $y(x) = (c_1 + c_2 x + c_3 x^2) e^x$.
Derivatives:
\[ y'(x) = (c_1 + c_2 + (c_2 + 2c_3)x + c_3 x^2) e^x \]
\[ y''(x) = (c_1 + 2c_2 + 2c_3 + (c_2 + 4c_3)x + c_3 x^2) e^x \]
Applying ICs at $x=0$:
\[ y(0) = c_1 = 1 \]
\[ y'(0) = c_1 + c_2 = 2 \implies 1 + c_2 = 2 \implies c_2 = 1 \]
\[ y''(0) = c_1 + 2c_2 + 2c_3 = 3 \implies 1 + 2(1) + 2c_3 = 3 \implies 2c_3 = 0 \implies c_3 = 0 \]
Particular Solution:
\[ y(x) = (1 + x) e^x \]

\subsubsection*{Solution 14: $y^{(4)} - y = 0, \, y(0)=1, \, y'(0)=0, \, y''(0)=1, \, y'''(0)=0$}
Roots: $r = \pm 1, \pm i$.
General solution: $y(x) = c_1 e^x + c_2 e^{-x} + c_3 \cos(x) + c_4 \sin(x)$.
Derivatives:
\[ y'(x) = c_1 e^x - c_2 e^{-x} - c_3 \sin(x) + c_4 \cos(x) \]
\[ y''(x) = c_1 e^x + c_2 e^{-x} - c_3 \cos(x) - c_4 \sin(x) \]
\[ y'''(x) = c_1 e^x - c_2 e^{-x} + c_3 \sin(x) - c_4 \cos(x) \]
Applying ICs at $x=0$:
\[ c_1 + c_2 + c_3 = 1 \]
\[ c_1 - c_2 + c_4 = 0 \]
\[ c_1 + c_2 - c_3 = 1 \]
\[ c_1 - c_2 - c_4 = 0 \]
From equations 2 and 4: $2c_4 = 0 \implies c_4 = 0$, so $c_1 = c_2$.
From equations 1 and 3: subtracting gives $2c_3 = 0 \implies c_3 = 0$.
So $c_1 + c_2 = 1 \implies 2c_1 = 1 \implies c_1 = 1/2, c_2 = 1/2$.
Particular Solution:
\[ y(x) = \frac{1}{2} e^x + \frac{1}{2} e^{-x} = \cosh(x) \]

\subsubsection*{Solution 15: $y''' + y'' + y' + y = 0, \, y(0)=2, \, y'(0)=-1, \, y''(0)=0$}
Roots: $r^2(r+1) + (r+1) = (r+1)(r^2+1) = 0 \implies r = -1, \pm i$.
General solution: $y(x) = c_1 e^{-x} + c_2 \cos(x) + c_3 \sin(x)$.
Derivatives:
\[ y'(x) = -c_1 e^{-x} - c_2 \sin(x) + c_3 \cos(x) \]
\[ y''(x) = c_1 e^{-x} - c_2 \cos(x) - c_3 \sin(x) \]
Applying ICs:
\[ c_1 + c_2 = 2 \]
\[ -c_1 + c_3 = -1 \]
\[ c_1 - c_2 = 0 \implies c_1 = c_2 \]
Since $c_1 = c_2$ and $c_1 + c_2 = 2 \implies c_1 = 1, c_2 = 1$.
Then $-1 + c_3 = -1 \implies c_3 = 0$.
Particular Solution:
\[ y(x) = e^{-x} + \cos(x) \]

\subsection*{Part C: Reconstructing Differential Equations}

\subsubsection*{Solution 16: $S = \{1, e^{2x}, e^{-2x}\}$}
Roots: $r_1 = 0, r_2 = 2, r_3 = -2$.
Polynomial: $P(r) = r(r-2)(r+2) = r(r^2 - 4) = r^3 - 4r = 0$.
Differential Equation:
\[ y''' - 4y' = 0 \]

\subsubsection*{Solution 17: $S = \{e^x, x e^x, x^2 e^x\}$}
Roots: $r = 1$ with multiplicity 3.
Polynomial: $P(r) = (r-1)^3 = r^3 - 3r^2 + 3r - 1 = 0$.
Differential Equation:
\[ y''' - 3y'' + 3y' - y = 0 \]

\subsubsection*{Solution 18: $S = \{\cos(3x), \sin(3x), e^x\}$}
Roots: $r_{1,2} = \pm 3i, r_3 = 1$.
Polynomial: $P(r) = (r^2 + 9)(r - 1) = r^3 - r^2 + 9r - 9 = 0$.
Differential Equation:
\[ y''' - y'' + 9y' - 9y = 0 \]

\subsubsection*{Solution 19: $S = \{1, x, e^{4x}\}$}
Roots: $r = 0$ (multiplicity 2), $r = 4$.
Polynomial: $P(r) = r^2(r - 4) = r^3 - 4r^2 = 0$.
Differential Equation:
\[ y''' - 4y'' = 0 \]

\subsubsection*{Solution 20: $S = \{e^{-x}, e^x \cos(2x), e^x \sin(2x)\}$}
Roots: $r_1 = -1$, $r_{2,3} = 1 \pm 2i$.
Factors: $(r+1)$ and $((r-1)^2 + 4) = (r^2 - 2r + 5)$.
Polynomial: $P(r) = (r+1)(r^2 - 2r + 5) = r^3 - 2r^2 + 5r + r^2 - 2r + 5 = r^3 - r^2 + 3r + 5 = 0$.
Differential Equation:
\[ y''' - y'' + 3y' + 5y = 0 \]

\subsection*{Part D: Logic and Clue Problems}

\subsubsection*{Solution 21: Complex Solution Clue}
\begin{enumerate}[label=(\alph*)]
    \item A solution of the form $y_1(x) = e^{3x}\sin(2x)$ implies a complex root $r_1 = 3 + 2i$. Since the differential equation has \textit{real coefficients}, complex roots must occur in conjugate pairs. Thus, $r_2 = 3 - 2i$ must also be a root.
    Therefore, the \textbf{minimal order is 2}.
    \item The characteristic factor is $(r - (3+2i))(r - (3-2i)) = (r-3)^2 + 4 = r^2 - 6r + 13 = 0$.
    The minimal differential equation is:
    \[ y'' - 6y' + 13y = 0 \]
    \item If $y_2(x) = 5 = 5e^{0x}$ is also a solution, then $r_3 = 0$ is a third root.
    The new \textbf{minimal order is 3}.
    Polynomial: $P(r) = r(r^2 - 6r + 13) = r^3 - 6r^2 + 13r = 0$.
    Differential Equation:
    \[ y''' - 6y'' + 13y' = 0 \]
\end{enumerate}

\subsubsection*{Solution 22: Repeated Complex Root Clue}
\begin{enumerate}[label=(\alph*)]
    \item The solution $y_1(x) = x e^{-x} \cos(x)$ has an explicit $x$ factor multiplying the exponential-trigonometric function. This implies that the complex conjugate pair $r = -1 \pm i$ ($\alpha = -1, \beta = 1$) has a \textbf{multiplicity of at least 2}.
    Roots: $r_1 = r_2 = -1 + i$, $r_3 = r_4 = -1 - i$.
    \item Since there are 2 pairs of complex roots, the \textbf{minimal order is $n = 4$}.
    \item The base quadratic factor is $(r - (-1+i))(r - (-1-i)) = (r+1)^2 + 1 = r^2 + 2r + 2$.
    Since the multiplicity is 2, the characteristic polynomial is:
    \[ P(r) = (r^2 + 2r + 2)^2 = r^4 + 4r^3 + 8r^2 + 8r + 4 = 0 \]
    The differential equation in operator notation is:
    \[ (D^2 + 2D + 2)^2 y = 0 \implies y^{(4)} + 4y''' + 8y'' + 8y' + 4y = 0 \]
\end{enumerate}

\subsubsection*{Solution 23: Reconstruction from Partial Solutions and IVP}
\begin{enumerate}[label=(\alph*)]
    \item The solution $y_1(x) = \sin(x)$ implies roots $r_{1,2} = \pm i$ ($\alpha = 0, \beta = 1$).
    The solution $y_2(x) = e^{2x}$ implies root $r_3 = 2$.
    Since the equation is of order 3, these 3 roots complete the characteristic spectrum: $\{i, -i, 2\}$.
    Polynomial: $P(r) = (r^2 + 1)(r - 2) = r^3 - 2r^2 + r - 2 = 0$.
    Differential Equation:
    \[ y''' - 2y'' + y' - 2y = 0 \]
    \item General solution: $y(x) = c_1 \cos(x) + c_2 \sin(x) + c_3 e^{2x}$.
    Derivatives:
    \[ y'(x) = -c_1 \sin(x) + c_2 \cos(x) + 2c_3 e^{2x} \]
    \[ y''(x) = -c_1 \cos(x) - c_2 \sin(x) + 4c_3 e^{2x} \]
    Applying ICs at $x=0$:
    \[ y(0) = c_1 + c_3 = 0 \implies c_1 = -c_3 \]
    \[ y'(0) = c_2 + 2c_3 = 3 \]
    \[ y''(0) = -c_1 + 4c_3 = 5 \implies -(-c_3) + 4c_3 = 5 \implies 5c_3 = 5 \implies c_3 = 1 \]
    Then $c_1 = -1$.
    From $c_2 + 2(1) = 3 \implies c_2 = 1$.
    Particular Solution:
    \[ y(x) = -\cos(x) + \sin(x) + e^{2x} \]
\end{enumerate}

\end{document}
